\begin{ledgroupsized}[r]{120mm}
\footnotesize
\pstart
\noindent\textbf{Überlieferung:}
\pend
\end{ledgroupsized}
\begin{ledgroupsized}[r]{114mm}
\footnotesize
\pstart
\parindent -6mm
\makebox[6mm][l]{\textit{L}}%
Notiz:
LH~XXXV~8, 30~Bl.~121. 
Ein Blatt~2\textsuperscript{o};
\href{\Scan{LH_35_08_30_121wz}}%
{Gegenmarke} eines Wasserzeichens in der Blattmitte (französisches Papier);
alle Ränder beschnitten.
Eineinviertel Seiten; 
Bl. 121~r\textsuperscript{o} horizontal in Felder aufgeteilt und vollständig beschrieben
mit Notizen meist mathematischen Inhalts (Kurven und Tangenten), 
Bl. 121~v\textsuperscript{o} bis auf zwei Bereiche, die Notizen auf Bl. 121~r\textsuperscript{o} in gleicher Höhe weiterführen, leer.
Ein Textfeld im unteren Drittel von Bl.~121~r\textsuperscript{o} überliefert den Anfang von N.\;\ref{Cc2_847-2} und ein Textbereich auf Bl.~121~v\textsuperscript{o} dessen Schlussteil.
Ein weiteres Textfeld im oberen Drittel von Bl.~121~r\textsuperscript{o} überliefert die Notiz "Calculus Elasticus et similes" (N.\;\ref{Cc2_847-1} in diesem Band).
\pend
\pstart
\noindent
Cc 2, Nr. 847
\pend
\end{ledgroupsized}
\vspace{5mm}
\begin{ledgroup}
\footnotesize
\pstart
\noindent
\textbf{Datierungsgründe:}
Mit der Handschrift LH~XXXV, 8, 30 Bl.~121 liegt
eine Sammlung sequentiell verfasster Notizen
zu den Themen Tangenten und "methodus tangentium inversa" vor,
die Leibniz ursprünglich wohl zu Streifen zerschneiden wollte. 
Deshalb hat er bei N.\;\ref{Cc2_847-2} und einer anderen Notiz, deren Länge eine Fortsetzung auf der Rückseite erforderte, den Schlussteil in gleicher Höhe wie den Anfang geschrieben. (Eine ähnliche Textverteilung durch sequentielle Abfassung liegt bei den Aufzeichnungen N.\;\ref{Cc2-1421} von 1679--1687 vor.).
\pend
\pstart
Im Zusammenhang mit dem allen gemeinsamen Thema Tangenten behandeln zwei dieser Notizen Fragen der Elastizität ("Calculus Elasticus et similes", N.\;\ref{Cc2_847-1}) bzw. der Statik ("De curvis inveniendis ope vectium conjugatorum", N.\;\ref{Cc2_847-2}). Sie werden daher in der Achten Reihe ediert.
\pend
\pstart
Die Datierung dieser Stücke beruht auf dem Wasserzeichen sowie auf dem engen Zusammenhang mit jeweils einem Komplex bereits edierter Texte aus dem Jahr 1674.
\pend
\pstart
Die Notiz N.\;\ref{Cc2_847-1} "Calculus elasticus et similes" wurde im oberen Bereich von Bl.~121~r\textsuperscript{o}, somit verhältnismäßig früh (und sicher früher als die Notiz N.\;\ref{Cc2_847-2}).
In N.\;\ref{Cc2_847-1} bezieht sich Leibniz mit den Worten "Videntur haec non satis consentire alibi dictis" auf eine eng zusammenhängende Gruppe eigenhändig auf Dezember 1674 datierter, den selben Fragen gewidmeter Texte: \cite{02608}\textit{Schediasma de calculo elastico} (\textit{LSB} VIII,~1 N.\;53, S.\;505--526) und \cite{02609}\textit{Calculus Elasticus ubi et quomodo alia problemata methodi tang. inversae reducantur subinde ad quadraturas} (ebd. N.\;54, S.\;527--540).
Zu dieser Gruppe ist außerdem das mathematische Konzept
\cite{00261}\textit{De progressionibus et geometria arcana et methodo tangentium inversa} (ebenfalls im "Xb.\;1674" verfasst; \textit{LSB} VII,~3 N.\;39) zu rechnen, das die theoretischen Grundlagen der anderen Texte erforscht; dort hat Leibniz in einem Zusatz bemerkt "Adde \textit{Schediasma de calculo elastico} \bracks{...} eodem mense".
\pend
\pstart
Da die auf Bl.~121 im Zusammenhang mit Fragen der Tangentenbestimmung und -methode entstandene Notiz N.\;\ref{Cc2_847-1} einerseits eng mit den Konzepten von Dezember 1674 verwandt ist, andererseits sich punktuell davon zu distanzieren scheint, kann in inhaltlicher Hinsicht von einer relativ zeitnahen Entstehung von N.\;\ref{Cc2_847-1} zu den obengenannten Konzepten ausgegangen werden.
\pend
\pstart
Im Einklang mit dieser Feststellung stehen auch die Befunde zum Wasserzeichen und zur Notiz "De curvis inveniendis", N.\;\ref{Cc2_847-2}. 
\pend
\pstart
Das Wasserzeichen in Bl.~121 ist für die zweite Hälfte von 1674 öfter belegt. Insbesondere lässt es sich auch im Träger des Stücks "De vectibus conjugatis 1" (\textit{LSB} VIII,~2 N.\;45\textsubscript{1}) feststellen, eines von vier Konzepten zu diesem Thema aus der zweiten Hälfte von 1674 (ebd. N.\;45).
Deshalb kann zunächst von einer Entstehung der Notizen auf Bl.~121 ab Mitte 1674 ausgegangen werden; wegen der sequentiellen Beschreibung des Blattes gilt außerdem der Terminus post quem von N.\;\ref{Cc2_847-1} (Dezember 1674) auch für alle nachfolgenden Notizen, darunter N.\;\ref{Cc2_847-2}.
\pend
\pstart
Leibniz beruft sich in N.\;\ref{Cc2_847-2} ausdrücklich, wenngleich inhaltlich vage, auf seine Texte über aufeinander wirkende Hebel ("vide schedas meas de vecte ducente"). Dabei handelt es sich höchstwahrscheinlich um die genannten Konzepte von 1674, insbesondere "De vectibus conjugatis 1" (VIII,~2 N.\;451, bes. S.\;391), das dasselbe Wasserzeichen wie der vorliegende Träger aufweist.
Dass Leibnizens vage Anspielung stattdessen den beiden Hannoveraner Konzepte "De vectibus conjugatis et figuris dentium in rotis" vom 25.\;Oktober (4.\;November) 1679 (N.\;\ref{intro_vectibus-conjugatis} in diesem Band) gelten könnte, kann zwar nicht zwingend ausgeschlossen werden, erscheint aber angesichts der weiteren aufgeführten Gründe als nicht plausibel. 
Auch im Fall von N.\;\ref{Cc2_847-2} kann daher von einer relativ zeitnahen Entstehung zur Textgruppe von 1674 ausgegangen werden.
\pend
\pstart
Zusammenfassend wird für die Notizen N.\;\ref{Cc2_847-1} und N.\;\ref{Cc2_847-2} die gemeinsame Datierung auf Dezember 1674 bis etwa 1675 vorgeschlagen, wobei N.\;\ref{Cc2_847-1} zuerst entstanden sein muss.
Eine spätere Entstehung, besonders der Notiz N.\;\ref{Cc2_847-2} im unteren Drittel des Blattes, kann nach heutigem Kenntnisstand nicht ausgeschlossen werden. Weitere Erkenntnisse sind aus der Herausgabe der übrigen Notizen auf dem gemeinsam Textträger (in einem künftigen Band der Siebenten Reihe) zu erwarten.
\pend
\end{ledgroup}
\selectlanguage{latin}		%
\vspace{8mm}
\pstart
\noindent
\href{\Scan{LH_35_08_30_121r}}%
{\lbrack121~r\textsuperscript{o}\rbrack} 
Invenire curvam
\edtext{\textit{BB}}{%
\lemma{}%
\Bfootnote{\textit{BB} \textit{erg.~L}}}
cujus ea sit natura ut sumto in ejus tangente \textit{BT} puncto \textit{E}, distantia ipsius \textit{E} a puncto dato
\edtext{\textit{D}.}{%
\lemma{}%
\Bfootnote{\textit{D}. \textit{erg.~L}}}
sit semper eadem data, 
\edtext{vide 
\edtext{schedas meas}{%
\lemma{schedas}%
\Bfootnote{%
\textit{(1)}~de
\textit{(2)}~meas~\textit{L}}}
de vecte ducente.\indexS{vectis ducens}}{%
\lemma{vide \bracks{...} ducente}%
\Cfootnote{Siehe die vier Konzepte \cite{02160}"De vectibus conjugatis" aus der Pariser Zeit (2.\;Hälfte 1674), \textit{LSB} VIII,\;2 N.\;45, bes. N.\;45\textsubscript{1}, S.\;391.}}
Venit nunc in mentem eam describi posse artificio corporis tracti;\indexS{corpus tractum}
\edtext{nempe radii mobilis}{%
\lemma{nempe}%
\Bfootnote{%
\textit{(1)}~radio mobili
\textit{(2)}~radii mobilis~\textit{L}}}
circa \textit{D} extremitati \textit{E} alligetur funiculus\indexS{funiculus}\indexS{funiculus secum trahens corpus grave} 
secum trahens corpus\indexS{corpus tractum} grave, 
quod planum in eodem semper puncto tangat. Hoc grave in plano\indexS{grave in plano} describet lineam quaesitam. Imo jam
\edtext{video lineae quaesitae exemplum unum esse circulum.}{%
\lemma{video}%
\Bfootnote{%
\textit{(1)}~lineam quaesitam non aliter esse quam Circulum
\textit{(2)}~lineae quaesitae exemplum unum esse circulum.~\textit{L}}}
\pend
\begin{edfigure}
\includegraphics[scale=1]{%
gesamttex/edit_VIII-4/images/LH_35_08_30_121_d1_121r.pdf}
\edcaption[.1em]{\lbrack\textit{Fig.~1}\rbrack}
\end{edfigure}
\pstart
Quaeritur curva 
\edtext{talis \textit{DE} figura}{%
\lemma{}%
\Bfootnote{%
talis \bars{\textit{DE} \textit{erg.}}
\textit{(1)}~ut
\textit{(2)}~figura~\textit{L}}}
scilicet dentis\indexS{dens rotae} a
\edtext{vecte\indexS{vectis} \textit{EF} (circa}{%
\lemma{vecte \textit{EF}}%
\Bfootnote{%
\textit{(1)}~ducendi
\textit{(2)}~(circa~\textit{L}}}
\textit{E} moto) ducendi, ut inter ducendum \textit{EF} semper tangat \textit{DF}
\edtext{in \textit{F}. Si eodem}{%
\lemma{}%
\Bfootnote{%
in \textit{F}. \bars{Fingamus oculum in \textit{E} positum. \textit{gestr.}} 
Si eodem~\textit{L}}}
manente dente quem fingimus jam constructum agatur \textit{E} in circulo circa \textit{D}, patet curvam fore illam ipsam quam paulo ante diximus quae filum\indexS{filum}
\edtext{\textit{EF}}{\lemma{}\Bfootnote{\textit{EF} \textit{erg.~L}}}
alligando radio \textit{DE}, in \textit{E}, describitur a fili extremo \textit{F}. Notabile hoc artificium per suppositum motum vel quietem, reddendi problema 
\edlabel{35_08_30_121_1a}%
\edtext{}{%
\xxref{35_08_30_121_1a}{35_08_30_121_1b}%
\lemma{}%
\Bfootnote{%
facilius. \bars{Alia \textit{gestr.}}
\href{\Scan{LH_35_08_30_121v}}%
{\bracks{121~v\textsuperscript{o}}} Quando~\textit{L}}}%
facilius.
\pend
\pstart
\href{\Scan{LH_35_08_30_121v}}%
{\lbrack121~v\textsuperscript{o}\rbrack} 
Quando\edlabel{35_08_30_121_1b} initio
\edtext{filum \textit{EF} alligatum}{%
\lemma{filum}%
\Bfootnote{%
\textit{(1)}~\textit{FE} i
\textit{(2)}~\textit{EF}
\textit{(a)}~sic collocatur, ut
\textit{(b)}~alligatum~\textit{L}}}
ad radium \textit{DE} in \textit{E}, sic collocatur ut angulus \textit{DEF} sit rectus inter gyrandum manebit rectus, et eo casu linea quaesita erit circulus. Verum eo casu non poterit dens\indexS{dens rotae} \textit{EF} motu suo circumagere dentem curvum\indexS{dens curvus} quia pellit ipsum versus centrum. Caeterum alia ratio se videtur offerre inveniendae hujus curvae. Si fingamus punctum \textit{E} immotum, et \textit{EF} agi circa \textit{E}, patet lineam dentis curvi\indexS{dens curvus} 
\edtext{\textit{DF},}{%
\lemma{}%
\Bfootnote{\textit{DF} \textit{erg.~L}}}
dum gyratur \textit{D} circa \textit{E}, semper circulum \textit{FF} secare ad angulos rectos. Sed nihil hinc.
\pend
\begin{edfigure}
\includegraphics[scale=1]{%
gesamttex/edit_VIII-4/images/LH_35_08_30_121_d2_121v.pdf}
\edcaption{\lbrack\textit{Fig.~2}\rbrack\hspace*{.25em}}
\end{edfigure}
